\uuid{PwkL}
\titre{Estimation des paramètres d'une loi normale}
\chapitre{Statistique}
\niveau{L2}
\module{Analyse de données}
\sousChapitre{Estimation}
\theme{statistiques}
\auteur{Maxime NGUYEN}
\organisation{AMSCC}
\datecreate{2024-09-01}
\difficulte{2}
\contenu{
\texte{$\mathcal N(\mu,\sigma)$ désigne la loi normale de moyenne $\mu$ et d’écart-type $\sigma$.}
\texte{Soient $X_1,\ldots,X_n$ indépendantes de loi $\mathcal N(\mu,\sigma)$, $\sigma>0$. On suppose $n\geq2$ et les observations $x_i$ non toutes égales. La densité est $f(x)=(\sigma\sqrt{2\pi})^{-1}\exp\left(-\frac{(x-\mu)^2}{2\sigma^2}\right)$.}
\begin{enumerate}
\item \question{Exprimer la vraisemblance et son logarithme.}
\reponse{$L(\mu,\sigma)=(\sigma\sqrt{2\pi})^{-n}\exp\left(-\frac{\sum(x_i-\mu)^2}{2\sigma^2}\right)$ et $\ell=-n\log\sigma-\frac n2\log(2\pi)-\frac{\sum(x_i-\mu)^2}{2\sigma^2}$.}
\item \question{Dériver $\ell$ par rapport à $\mu$.}
\reponse{$\frac{\partial\ell}{\partial\mu}=\frac{\sum(x_i-\mu)}{\sigma^2}$.}
\item \question{En déduire l'estimation et l'estimateur de $\mu$.}
\reponse{La dérivée s'annule en $\widehat\mu_{\mathrm{obs}}=\bar x$. C'est un maximum pour chaque $\sigma$, car $\frac{\partial^2\ell}{\partial\mu^2}=-\frac{n}{\sigma^2}<0$. L'estimateur est $\widehat M_n=\overline X$.}
\item \question{Déterminer de même l'estimation et l'estimateur de $\sigma$.}
\reponse{Après remplacement de $\mu$ par $\bar x$, $\frac{\partial\ell}{\partial\sigma}=-\frac{n}{\sigma}+\frac{\sum(x_i-\bar x)^2}{\sigma^3}$. Le maximum est donc
\[\widehat\sigma_{\mathrm{obs}}=\sqrt{\frac1n\sum_{i=1}^n(x_i-\bar x)^2},\qquad
\widehat H_n=\sqrt{\frac1n\sum_{i=1}^n(X_i-\overline X)^2}.\]
La décomposition $\sum(x_i-\mu)^2=\sum(x_i-\bar x)^2+n(\bar x-\mu)^2$ justifie le maximum conjoint. Le diviseur de l'EMV de variance est $n$, tandis que l'estimateur sans biais de variance utilise $n-1$.
Si toutes les observations étaient égales, la vraisemblance serait non bornée lorsque $\sigma\to0$ et il n'y aurait pas de maximum pour $\sigma>0$.}
\end{enumerate}
}
